\documentclass[11pt,letterpaper]{article}
\usepackage[utf8]{inputenc}
\usepackage[T1]{fontenc}
\usepackage[provide=*,spanish]{babel}
\usepackage[margin=2cm,headheight=14pt]{geometry}
\usepackage{lmodern,microtype,amsmath,amssymb,array,tabularx,longtable,booktabs,enumitem,fancyhdr,listings,upquote}
\usepackage[hidelinks]{hyperref}
\setlength{\parindent}{0pt}\setlength{\parskip}{5pt}
\setlength{\emergencystretch}{3em}
\setlist{leftmargin=*,itemsep=2pt,topsep=3pt}
\renewcommand{\arraystretch}{1.1}
\pagestyle{fancy}\fancyhf{}\fancyhead[L]{\small Team 07 -- Parser y SDT}\fancyhead[R]{\small Model Designer}
\fancyfoot[L]{\small Brandon Yáñez Barajas}\fancyfoot[R]{\thepage}
\lstset{basicstyle=\ttfamily\small,columns=fullflexible,keepspaces=true,breaklines=true,showstringspaces=false,upquote=true,frame=none}
\newcommand{\titulo}[1]{{\Large\bfseries #1}\par\vspace{5pt}}

\begin{document}
\titulo{1. Gramática y símbolos de entrada}
\textbf{Unidad de análisis:} \texttt{int main() Block}. Se conservan las producciones 0--50 y la numeración de los 161 estados CLR(1).
\section*{Terminales y normalización}
La entrada está formada por tokens del lexer anterior. Cada token conserva \texttt{type}, \texttt{lexeme}, \texttt{line} y \texttt{column}. La etiqueta gramatical normalizada es \texttt{terminal}. El parser trabaja con tokens, no con caracteres sueltos.

\begin{tabularx}{\linewidth}{@{}p{4cm}X@{}}\toprule Entrada del lexer & Terminal del parser\\\midrule
Palabra clave admitida & Su lexema: \texttt{int, float, char, if, else, while, return}.\\
Identificador & \texttt{ID}; el lexema \texttt{main} se normaliza como \texttt{main}.\\
Operador o puntuación & Su lexema completo: \texttt{==} no se forma uniendo dos tokens \texttt{=}.\\
Constante válida & \texttt{INT, FLOAT, CHAR, STRING}, según su clase léxica.\\
Fin de entrada & \texttt{\$}, una sola vez. Un error léxico bloquea el parser.\\\bottomrule\end{tabularx}

Los 31 terminales del lenguaje son:
\texttt{! != \% ( ) * + - / ; \textless{} \textless{}= = == \textgreater{} \textgreater{}= CHAR FLOAT ID INT STRING char else float if int main return while \{ \}}.
\texttt{int} es una palabra clave; \texttt{INT}, una constante. \texttt{epsilon} no es token. ACTION añade \texttt{\$} como columna 32.
\section*{No terminales y símbolo inicial}
\texttt{Program, Block, Items, Item, Decl, Type, Init, Stmt, Matched, Unmatched, Simple, Expr, Equality, Relation, Add, Mul, Unary, Primary}.
El símbolo inicial es \texttt{Program}. La aumentación agrega \texttt{Start -> Program} (p0). Las listas vacías y los inicializadores ausentes tienen una parte derecha vacía.
\section*{Unidad de entrada y declaraciones}
La unidad completa reconocida por Program contiene una función \texttt{int main()} y su bloque. Por ejemplo:
\begin{lstlisting}[language=C]
int main() {
    int x = 10;
    return x;
}
\end{lstlisting}
La declaración \texttt{int x = 10;} es válida dentro del bloque, pero no constituye un Program completo por sí sola. Las declaraciones y las asignaciones se distinguen por su estructura: la declaración comienza por Type y la asignación por ID.
\clearpage\section*{Producciones completas}
La numeración es la usada por las reducciones y las acciones SDT.
{\footnotesize\renewcommand{\arraystretch}{1.0}\begin{longtable}{@{}r p{10.2cm} p{4.6cm}@{}}\toprule ID & Producción & Función\\\midrule\endhead
0 & \texttt{Start -\textgreater{} Program} & Regla aumentada; aceptación.\\
1 & \texttt{Program -\textgreater{} int main ( ) Block} & Unidad de entrada.\\
2 & \texttt{Block -\textgreater{} \{ Items \}} & Bloque y ámbito.\\
3 & \texttt{Items -\textgreater{} Items Item} & Agrega un elemento.\\
4 & \texttt{Items -\textgreater{} epsilon} & Lista vacía.\\
5 & \texttt{Item -\textgreater{} Decl} & Elemento de declaración.\\
6 & \texttt{Item -\textgreater{} Stmt} & Elemento de sentencia.\\
7 & \texttt{Decl -\textgreater{} Type ID Init ;} & Variable e inicializador opcional.\\
8 & \texttt{Type -\textgreater{} int} & Tipo entero.\\
9 & \texttt{Type -\textgreater{} float} & Tipo real.\\
10 & \texttt{Type -\textgreater{} char} & Tipo carácter.\\
11 & \texttt{Init -\textgreater{} = Expr} & Inicializador presente.\\
12 & \texttt{Init -\textgreater{} epsilon} & Sin inicializador.\\
13 & \texttt{Stmt -\textgreater{} Matched} & Sentencia cerrada.\\
14 & \texttt{Stmt -\textgreater{} Unmatched} & Sentencia abierta.\\
15 & \texttt{Matched -\textgreater{} Simple} & Instrucción simple.\\
16 & \texttt{Matched -\textgreater{} Block} & Bloque delimitado.\\
17 & \texttt{Matched -\textgreater{} if ( Expr ) Matched else Matched} & if/else cerrado.\\
18 & \texttt{Matched -\textgreater{} while ( Expr ) Matched} & while de cuerpo cerrado.\\
19 & \texttt{Unmatched -\textgreater{} if ( Expr ) Stmt} & if pendiente de else.\\
20 & \texttt{Unmatched -\textgreater{} if ( Expr ) Matched else Unmatched} & if/else de cierre pendiente.\\
21 & \texttt{Unmatched -\textgreater{} while ( Expr ) Unmatched} & while de cuerpo abierto.\\
22 & \texttt{Simple -\textgreater{} ID = Expr ;} & Asignación.\\
23 & \texttt{Simple -\textgreater{} return Expr ;} & Retorno con expresión.\\
24 & \texttt{Simple -\textgreater{} ;} & Sentencia vacía.\\
25 & \texttt{Expr -\textgreater{} Equality} & Entrada a expresión.\\
26 & \texttt{Equality -\textgreater{} Equality == Relation} & Igualdad.\\
27 & \texttt{Equality -\textgreater{} Equality != Relation} & Desigualdad.\\
28 & \texttt{Equality -\textgreater{} Relation} & Paso a relaciones.\\
29 & \texttt{Relation -\textgreater{} Relation \textless{} Add} & Menor.\\
30 & \texttt{Relation -\textgreater{} Relation \textless{}= Add} & Menor o igual.\\
31 & \texttt{Relation -\textgreater{} Relation \textgreater{} Add} & Mayor.\\
32 & \texttt{Relation -\textgreater{} Relation \textgreater{}= Add} & Mayor o igual.\\
33 & \texttt{Relation -\textgreater{} Add} & Paso a suma.\\
34 & \texttt{Add -\textgreater{} Add + Mul} & Suma.\\
35 & \texttt{Add -\textgreater{} Add - Mul} & Resta.\\
36 & \texttt{Add -\textgreater{} Mul} & Paso a producto.\\
37 & \texttt{Mul -\textgreater{} Mul * Unary} & Multiplicación.\\
38 & \texttt{Mul -\textgreater{} Mul / Unary} & División.\\
39 & \texttt{Mul -\textgreater{} Mul \% Unary} & Módulo entero.\\
40 & \texttt{Mul -\textgreater{} Unary} & Paso a unario.\\
41 & \texttt{Unary -\textgreater{} + Unary} & Positivo unario.\\
42 & \texttt{Unary -\textgreater{} - Unary} & Negativo unario.\\
43 & \texttt{Unary -\textgreater{} ! Unary} & Negación lógica.\\
44 & \texttt{Unary -\textgreater{} Primary} & Paso a primario.\\
45 & \texttt{Primary -\textgreater{} ID} & Uso de identificador.\\
46 & \texttt{Primary -\textgreater{} INT} & Constante entera.\\
47 & \texttt{Primary -\textgreater{} FLOAT} & Constante real.\\
48 & \texttt{Primary -\textgreater{} CHAR} & Literal carácter.\\
49 & \texttt{Primary -\textgreater{} STRING} & Literal cadena.\\
50 & \texttt{Primary -\textgreater{} ( Expr )} & Agrupación explícita.\\
\bottomrule\end{longtable}}
\clearpage\section*{Precedencia y asociatividad}
De mayor a menor prioridad: primarios y paréntesis; prefijos \texttt{+ - !}; \texttt{* / \%}; \texttt{+ -}; \texttt{< <= > >=}; \texttt{== !=}. Los prefijos se asocian a la derecha y los binarios a la izquierda. La asignación solo es sentencia o inicialización, no parte de Expr.

\texttt{Matched/Unmatched} asocia cada \texttt{else} con el \texttt{if} abierto más cercano. Las variantes de \texttt{while} conservan esa distinción. Las declaraciones como cuerpo de control requieren llaves.

\section*{Cómo reconocer una declaración}
Para \texttt{int x = 10;} dentro de main, el lexer entrega \texttt{int ID = INT ;}. Los lexemas \texttt{x} y \texttt{10} permanecen asociados a sus tokens. La gramática organiza esa secuencia con estas reglas:
\begin{lstlisting}
p7:  Decl -> Type ID Init ;
p8:  Type -> int
p11: Init -> = Expr
p46: Primary -> INT
\end{lstlisting}
Entre Expr y Primary se aplican las producciones de paso 25, 28, 33, 36, 40 y 44. No se omiten al construir el árbol. Así, la declaración tiene cuatro hijos: Type, ID, Init y el punto y coma.

En \texttt{int x;} se usa p12, \texttt{Init -> epsilon}. Init sigue siendo un nodo del árbol, pero no tiene hijos. La declaración no cambia de producción: sigue usando p7.
\clearpage\section*{Función de los grupos de no terminales}
\begin{tabularx}{\linewidth}{@{}p{4.3cm}X@{}}\toprule Grupo & Estructura que representa\\\midrule
Program y Block & La unidad completa y un cuerpo delimitado por llaves.\\
Items e Item & Una lista de declaraciones y sentencias; Items también permite una lista vacía.\\
Decl, Type e Init & Nombre, tipo e inicialización opcional de una variable.\\
Stmt, Matched y Unmatched & Sentencias y asociación no ambigua de if/else.\\
Simple & Asignación, retorno o punto y coma aislado.\\
Expr hasta Primary & Niveles de prioridad de las expresiones.\\\bottomrule\end{tabularx}
\section*{Ejemplos de prioridad y agrupación}
\texttt{2 + 3 * 4} se reconoce como \texttt{2 + (3 * 4)} porque Add utiliza un Mul completo como operando. \texttt{8 - 3 - 1} se agrupa como \texttt{(8 - 3) - 1} por la recursión izquierda de Add. Los paréntesis de \texttt{(2 + 3) * 4} fuerzan primero la suma.

La forma \texttt{if (x) if (y) ; else ;} vincula el else con el if de y. Matched representa la sentencia interior con else; Unmatched permite que el if exterior siga abierto. No se resuelve este caso dando prioridad manual al token else en ACTION.

\section*{Límites y verificación}
Se admiten bloques, declaraciones de una variable, asignaciones, expresiones, \texttt{if/else}, \texttt{while} y \texttt{return Expr;}. Se excluyen otras funciones y llamadas, declaraciones globales, arreglos, apuntadores, \texttt{for}, \texttt{switch}, \texttt{struct}, preprocesador y operadores no enumerados. \texttt{STRING} permite reconocer sintácticamente \texttt{int x="hello";} y rechazarlo después en SDT; no añade un tipo declarable de cadena.

La gramática produce 161 estados canónicos y 730 transiciones. Su tabla ACTION/GOTO no presenta conflictos shift/reduce ni reduce/reduce.
\end{document}
